\section{The banditGPT Framework}
\label{sec:methodology}

Building on the formulation and feature space established in Section~\ref{sec:preliminaries}, we now describe the banditGPT architecture. The system combines three core mechanisms to address the production challenges identified in Section~\ref{sec:intro}---preference heterogeneity, distribution drift, model non-stationarity, and the absence of hedging in existing bandit routers:

\begin{itemize}[nosep,leftmargin=*]
    \item \textbf{Cost-aware contextual bandits} initialized from offline warmup priors, using per-prompt embeddings for continuous online learning.
    \item \textbf{Corralling meta-learning} that hedges against prior mismatch via complementary exploration schedules, providing a safety mechanism absent from existing bandit routers.
    \item \textbf{Family-level parameter sharing} via Hybrid LinUCB, which pools observations across models within the same provider family to improve sample efficiency in larger portfolios.
\end{itemize}

\subsection{System Overview}
\label{sec:architecture}

The system consists of three distinct layers of decision-making, as illustrated in Figure~\ref{fig:architecture}.

\begin{figure}[t]
\centering
\includegraphics[width=\columnwidth]{figures/figure2_architecture.png}
\caption{\textbf{banditGPT architecture (overview).} A user prompt is encoded via Sentence-BERT and PCA into context $x_t \in \mathbb{R}^{33}$. The constraint filter prunes models by cost, latency, and quality floors to yield candidate set $\mathcal{A}_t$. A Corralling coordinator (Log-Barrier OMD; Agarwal et al., 2017) stochastically mixes two experts---a Hybrid LinUCB expert warm-started from offline priors and a Disjoint LinUCB tabula rasa expert---then the cost-aware UCB objective selects model $a_t$. Online reward $r_t$ feeds back to update both coordinator weights (via OMD) and expert parameters (via importance-weighted LinUCB updates). LLM portfolio shown is illustrative.}
\Description{A vertical pipeline routing diagram: prompt embedding and PCA produce context features; a constraint filter prunes models using a model registry; a Corralling coordinator (Log-Barrier OMD) mixes a Hybrid LinUCB warmup expert and a Disjoint LinUCB tabula rasa expert with probabilities $p_t(1)$ and $p_t(2)$; a cost-aware Hybrid LinUCB layer scores and selects a model from the LLM portfolio; observed reward feeds back to update coordinator weights and expert parameters.}
\label{fig:architecture}
\end{figure}

\subsection{Layer 1: Budget Constraint Filtering}
\label{sec:layer1}

To enforce operational constraints, we implement a per-request eligibility check before the bandit layer. For each candidate model $m \in \mathcal{M}_t$, the router applies hard filters:
\begin{equation}
\mathcal{E}_{x_t} = \{m \in \mathcal{M}_t \mid c_m \leq c_{\max},\; \ell_m \leq \ell_{\max},\; q_m \geq q_{\min}\}
\end{equation}
where $c_{\max}$ and $q_{\min}$ are user-specified cost ceiling and quality floor constraints respectively. Models failing any constraint are excluded before the bandit selects, guaranteeing per-request constraint satisfaction. Latency filtering is supported in the implementation but not evaluated in this paper. Cost and quality trade-offs among eligible models are then handled by the cost-penalized LinUCB objective (Layer 3).

\subsection{Layer 2: Expert Corralling for Adaptive Meta-Learning}
\label{sec:layer2}

Non-stationarity (e.g., a ``good'' model becoming ``bad'' due to drift) and misspecified offline priors pose significant challenges for standard bandit algorithms relying on fixed initializations. We address this \emph{at the meta level} via a strategy inspired by Expert Corralling~\cite{agarwal2017corralling}, which manages a portfolio of expert policies: a ``Warmup Expert'' initialized with static priors (constant $\alpha{=}2.0$) and a ``Tabula Rasa Expert'' with high plasticity (decaying $\alpha{=}1.0$ to $0.01$). Corralling exploits the \emph{complementarity} of these schedules: the warmup expert provides informed early selections, while the tabula rasa expert converges to the best policy over time. The meta-learner dynamically shifts weight between them, achieving 7 to 13\% lower regret than warmup-only across all corruption levels and up to ${\sim}6\%$ lower than tabula rasa at $\alpha{=}0$ (tapering as corruption increases), while enabling implicit avoidance of degraded models through per-arm UCB updates even when priors are adversarial (Section~\ref{sec:results}). Importantly, each expert bandit is itself trained under a \emph{stationary-reward assumption} (monotone information accumulation without explicit forgetting); prior mismatch is handled solely through the meta-learner's ability to reweight experts whose performance has degraded.

The master algorithm updates the probability distribution over experts $w_{i,t}$ using an exponential weight update based on importance-weighted loss estimates $\hat{\ell}_{i,t}$. To prevent ``Expert Death'' (where a temporarily poor expert's weight vanishes, rendering it unrecoverable), we implement a mixing parameter $\gamma$:
\begin{equation}
p_{i,t} = (1 - \gamma) \frac{w_{i,t}}{\sum_j w_{j,t}} + \frac{\gamma}{K}
\end{equation}

We set $\gamma=0.05$, imposing a uniform exploration floor that enables asymptotic recovery if the optimal expert changes over time.

\subsection{Layer 3: Cost-Aware Hybrid LinUCB}
\label{sec:layer3}

Each expert $i$ models the reward function using LinUCB~\cite{li2010contextual}, which maintains a regularized least-squares estimate of $\theta_m$ and an associated confidence set. The policy selects the arm maximizing the Upper Confidence Bound (UCB) of the composite utility:
\begin{equation}
a_t = \arg\max_{a \in \mathcal{E}_{x_t}} \left( x_t^\top \hat{\thetavec}_a + \alpha \sqrt{x_t^\top \Amat_a^{-1} x_t} - \lambda c_a \right)
\end{equation}

\paragraph{Family-Level Parameter Sharing.}
When the model portfolio contains models from the same provider family (e.g., multiple iterations of the same foundation model), we extend LinUCB with a \emph{Hybrid} formulation~\cite{li2010contextual} that decomposes the reward coefficient into a shared family component $\betavec_f$ and a per-model residual $\thetavec_m$:
\begin{equation}
\mathbb{E}[r | x, m] = x^\top (\betavec_{f(m)} + \thetavec_m)
\end{equation}
where $f(m)$ maps model $m$ to its family. The shared $\betavec_f$ is updated by \emph{every} observation from any family member, while $\thetavec_m$ captures model-specific deviations. When all families are singletons, the hybrid formulation reduces to disjoint LinUCB with independent per-model parameters. The family structure is operationally relevant during model onboarding (Section~\ref{sec:transfer}): when a new model joins an existing family, it immediately inherits the family's accumulated $\betavec$, providing a continuously updated reward estimate. A leave-one-out ablation (Section~\ref{sec:transfer_null}) confirms that this continuous family sharing subsumes one-shot per-model $\thetavec$ transfer, suggesting that the hybrid decomposition is an effective abstraction level for cross-model knowledge transfer in our evaluation.

\paragraph{Family Assignment.}
The mapping $f(m)$ can be supplied manually, extracted from metadata, or inferred syntactically (e.g., stripping suffixes like \texttt{gpt-4-turbo} $\to$ \texttt{gpt-4}). In practice, we provide a \texttt{compute\_correlation\_families()} utility that clusters models via connected components on within-provider tetrachoric correlations~\cite{drasgow1986polychoric}. This data-driven approach is used for our $K{\geq}5$ evaluations (Section~\ref{sec:results}).

\subsection{Meta-learning, Dynamic Portfolios, and Calibration}
\label{sec:transfer}
\label{sec:corralling}
Because LLM model availability and APIs evolve constantly, banditGPT must handle dynamic portfolios and potentially miscalibrated priors. 
We address this via the Corralling coordinator, which maintains non-zero probability on each expert to avoid starvation and updates weights based on observed rewards. As a precaution against potential \emph{Scale Explosion} when loading pre-trained priors with mismatched reward scales, we apply automatic two-pass calibration at initialization (bias clamping and suite probing). We defer the full meta-learning update rules, regularization schemes, and model onboarding mechanisms to the supplementary material.

